We record the geometric shadow of an extension-tiler factorization.  This
section uses only the standard order polytope of a poset and the standard
simplex triangulation of the cube; no projected or visual model is part of
the proof.

Let $P$ be a poset on a finite label set $X$, with $|X|=n$.  Its order
polytope is
\[
        \mathcal O(P)=
        \{x\in [0,1]^X : x_i\le x_j \text{ whenever } i<_P j\},
\]
in the sense of Stanley~\cite{Stanley1986}.  For a word
$\sigma=(\sigma_1,\ldots,\sigma_n)\in S_X$, let
\[
        \Delta_\sigma =
        \{x\in [0,1]^X :
        0\le x_{\sigma_1}\le x_{\sigma_2}\le\cdots
        \le x_{\sigma_n}\le 1\}.
\]
The closed simplices $\Delta_\sigma$, for $\sigma\in S_X$, form the
standard triangulation of the unit cube, with disjoint relative interiors.

\begin{lemma}
\label{lem:order-polytope-standard-triangulation}
For every finite poset $P$ on $X$,
\[
        \mathcal O(P)=\bigcup_{\sigma\in L(P)}\Delta_\sigma .
\]
\end{lemma}

\begin{proof}
If $\sigma$ is a linear extension of $P$, then $i<_P j$ implies that
$i$ occurs before $j$ in $\sigma$, so every point of $\Delta_\sigma$
satisfies $x_i\le x_j$.  Thus $\Delta_\sigma\subseteq \mathcal O(P)$.

Conversely, every point $x\in [0,1]^X$ lies in at least one simplex of the
standard triangulation, obtained by ordering the coordinates and breaking
ties arbitrarily.  If $x\in\mathcal O(P)$, the tie-breaking can be chosen
so that $i<_P j$ places $i$ no later than $j$.  The resulting word is a
linear extension of $P$, and $x$ lies in the corresponding
$\Delta_\sigma$.
\end{proof}

For $a\in S_X$, let $T_a$ be the coordinate permutation defined by
\[
        (T_a x)_{a(i)}=x_i
        \qquad (i\in X).
\]
Then $T_a(\Delta_\sigma)=\Delta_{a\sigma}$, where $a$ acts on the labels
in the word.  We write
\[
        a\mathcal O(P):=T_a(\mathcal O(P))
        =\bigcup_{\sigma\in L(P)}\Delta_{a\sigma},
\]
for this coordinate-relabeling of $\mathcal O(P)$ by $a$.

\begin{proposition}[Extension tilers give cube tilings]
\label{prop:extension-tiler-order-polytope-cube-tiling}
Suppose that $P$ is an extension tiler on $X$, so that
\[
        S_X=\bigsqcup_{a\in A} aL(P).
\]
Then the coordinate-permuted order polytopes $a\mathcal O(P)$ tile the
unit cube:
\[
        [0,1]^X=\bigcup_{a\in A} a\mathcal O(P),
\]
with overlaps only along faces of the standard simplex triangulation.
\end{proposition}

\begin{proof}
By Lemma~\ref{lem:order-polytope-standard-triangulation}, the pieces
$a\mathcal O(P)$ are unions of the standard cube simplices indexed by
the words in $aL(P)$.  The assumed disjoint union
$S_X=\bigsqcup_{a\in A}aL(P)$ assigns each standard simplex to exactly
one piece.  Since the standard simplices tile the cube and have disjoint
relative interiors, the stated cube tiling follows.
\end{proof}

\begin{corollary}[The six-element counterexample order-polytope cube tiling]
\label{cor:p038-order-polytope-cube-tiling}
For the poset $P$ of Theorem~\ref{thm:s6-counterexample}, called $P038$ in
the supplementary enumeration, the unit cube $[0,1]^6$ is tiled, up to
shared triangulation faces, by the $48$ coordinate permutations
$a\mathcal O(P)$ with $a\in A=HK\sqcup H\tau K$.  Each copy has volume
$15/720=1/48$.
\end{corollary}

\begin{proof}
Theorem~\ref{thm:s6-counterexample} gives
$S_6=\bigsqcup_{a\in A}aL(P)$ with $|A|=48$ and $|L(P)|=15$.
Apply Proposition~\ref{prop:extension-tiler-order-polytope-cube-tiling}.
The volume statement follows either from the $15$ standard simplices in
$\mathcal O(P)$ or from the $48$-piece cube tiling.
\end{proof}

\subsection{Certified geometry of the P038 tile}

The tiling proof above uses only the chamber partition.  We also record a
small amount of certified geometry for the single tile \(\mathcal O(P038)\),
so that the shape used in the corollary is independently inspectable.
The standard order-polytope coordinates in this section satisfy
\(x_i\le x_j\) when \(i<_P j\).  For the finite face data it is convenient to
use order-ideal indicator coordinates \(y_i=1-x_i\), in which a cover
\(i<_P j\) gives \(y_j\le y_i\).  This integral affine change of coordinates
does not change the face lattice, lattice-normalized volume, or Ehrhart data.

For the cover relations
\[
        0<1,\qquad 1<2,\qquad 1<3,\qquad
        4<3,\qquad 4<5,\qquad 5<2,
\]
the irredundant natural support inequalities in the \(y\)-coordinates are
\[
\begin{array}{lll}
  0\le y_2, & 0\le y_3, & y_0\le 1,\\
  y_4\le 1, & y_1\le y_0, & y_2\le y_1,\\
  y_3\le y_1, & y_3\le y_4, & y_5\le y_4,\\
  y_2\le y_5. &&
\end{array}
\]

\begin{proposition}[Certified \(P038\) order-polytope data]
\label{prop:p038-certified-order-polytope-data}
A finite replay verifies that \(\mathcal O(P038)\) has \(13\) vertices,
\(10\) irredundant natural support facets, and boundary \(f\)-vector
\[
        (f_0,f_1,f_2,f_3,f_4,f_5)
        =(13,50,88,81,40,10).
\]
The top face count is \(f_6=1\).  Its Ehrhart \(h^*\)-polynomial is
\[
        h^*_{\mathcal O(P038)}(t)=1+6t+7t^2+t^3,
\]
equivalently
\[
        L_{\mathcal O(P038)}(m)=
        \binom{m+6}{6}
        +6\binom{m+5}{6}
        +7\binom{m+4}{6}
        +\binom{m+3}{6}.
\]
\end{proposition}

\begin{proof}
The script \texttt{scripts/verify\_p038\_order\_polytope\_geometry.py}
enumerates the \(13\) order ideals, forms their \(0/1\) vertices, checks the
ten support facets above, enumerates all nonempty intersections of those
facet hyperplanes, and computes the affine dimension of each resulting face.
It also counts order-preserving maps \(P038\to\{0,\ldots,m\}\) for
\(m=0,\ldots,6\), derives the \(h^*\)-coefficients, and checks the displayed
binomial-basis Ehrhart formula.  These computations are supporting tile data;
the cube tiling itself is still proved by the chamber partition in
Proposition~\ref{prop:extension-tiler-order-polytope-cube-tiling}.
\end{proof}

\begin{remark}
This is a coordinate-permutation tiling inherited from the braid-arrangement
triangulation.  It is not a statement about arbitrary Euclidean congruent
copies, and it is independent of any three-dimensional projected model used
for visualization.
\end{remark}
